% !TeX root = ../main.tex

\chapter{相关工作}
\label{chap:related}

大规模图的在线推理效率是图神经网络面临的持续挑战。$L$层消息传递GNN的有效感受野在最坏情况下可达$O(\bar{d}^L)$，邻居收集与内存访问开销随层数快速放大。稀疏分解图神经网络（SDGNN）将多层GNN输出的逼近转化为单次稀疏矩阵乘，把推理复杂度降至$O(n\bar{k}d)$\cite{Hu2024SDGNN}；然而，其全局统一的Lasso正则系数$\lambda_{\mathrm{reg}}$对所有节点施加相同稀疏惩罚，无法适配不同拓扑角色节点在频域参与方式上的差异，是本文方法链的核心改进出发点。本章围绕三条文献脉络展开：第~\ref{sec:ch2-gnn-basics}~节建立全篇复用的基础符号体系与谱图理论背景；第~\ref{sec:ch2-efficient}~节梳理从采样加速、线性化预计算到SDGNN稀疏分解的技术演进，重点揭示其结构性局限；第~\ref{sec:ch2-heterophily}~节讨论复杂图结构上的频率响应与相关方法，为第4章频率接口设计提供背景动机。第~\ref{sec:ch2-summary}~节提炼三大研究空白并与第~3~章至第~5~章方法链逐一衔接。本章叙述范围限于与谱驱动自适应稀疏分解直接相关的背景，不引入本文原创定理、命题或算法，相关正式内容统一留待后续方法章展开。

% ---- 图~\ref{fig:ch2-roadmap}~占位符 ----
\begin{figure}[htbp]
  \centering
  \includegraphics[width=1\linewidth]{2-1.pdf}
      % \caption*{注：本图反映三条文献脉络对后续方法章节的知识铺垫关系，而非章节引用的严格对应。各脉络详细综述见正文第2.1--2.3节；研究空白与章节衔接见第2.4节。}
      % \setlength{\fboxsep}{1.2em}
      % \fbox{\parbox{0.88\textwidth}{\small\centering
      %   \textbf{【图2.1 占位符】}\\[0.6em]
      %   图题：本章三条文献脉络与后续方法章节的对应关系\\[0.4em]
      %   表达内容：图左侧纵向排列三条文献脉络“图神经网络基础与谱图理论—高效与稀疏图神经网络方法—复杂图结构与谱视角相关方法”，右侧纵向排列第~3~章至第~5~章的方法链“问题形式化与图复杂度指标体系—谱驱动三维自适应稀疏机制—算法整合、理论保证与复杂度分析”。箭头仅表示知识铺垫关系，其中前两条脉络主要支撑第~3~章与第~5~章的问题背景与理论闭环，第三条脉络主要支撑第~4~章的频率维设计。图注需注明：本图为章节衔接示意，不是算法流程图。
      % }}
      \caption{本章三条文献脉络与后续方法章节的对应关系}
  \label{fig:ch2-roadmap}
\end{figure}


%%============================================================
\section{图神经网络基础与谱图理论}
\label{sec:ch2-gnn-basics}

本节补充后续分析所需的核心符号，并引入谱图理论背景。全文统一符号见《符号和缩略语说明》，本章不再重复列示；本章新增符号均在首次出现处行内定义，涉及第~3~章至第~5~章方法接口的记号再于相应章节引言中集中说明。

无向图$G=(V,E,X)$由节点集$V$（$|V|=n$）、边集$E$和节点特征矩阵$X\in\mathbb{R}^{n\times D}$组成；邻接矩阵$A$为对称矩阵，度矩阵$D_\text{deg}$的对角元素$D_{\text{deg},ii}=d_i=\sum_j A_{ij}$，全图平均度$\bar{d}=m/n$，其中$m:=\mathrm{nnz}(A)=2|E|$。本文统一以$m$作为复杂度分析的图规模参数。归一化拉普拉斯矩阵定义为\label{idx:theme:normlaplacian}
\begin{equation}
  \tilde{\mathcal{L}} = I - D_\text{deg}^{-1/2}AD_\text{deg}^{-1/2},
  \label{eq:ch2-normlaplacian}
\end{equation}

其谱分解$\tilde{\mathcal{L}}=U\Lambda U^\top$中特征值满足$0=\lambda_0\leq\lambda_1\leq\cdots\leq\lambda_{n-1}\leq 2$。本文统一采用从 $\lambda_0$ 开始的编号方式；因此对连通图，谱间隙记为$\lambda_\text{gap}:=\lambda_1$，即最小非零特征值（也就是通常所说的第二小特征值）。若图不连通，则更一般地取最小正特征值。该量刻画图的连通扩张性，与导通率（conductance）之间满足Cheeger型关系\cite{Shuman2013GSP}，在后续跳距预算衰减分析中作为图级结构量使用。多项式谱滤波器$g(\tilde{\mathcal{L}})=\sum_{k=0}^{K_\text{poly}}c_k\tilde{\mathcal{L}}^k$可在顶点域以$O(K_\text{poly}\cdot m\cdot D)$计算，无需显式特征分解；对大规模图，Lanczos迭代以$O(m\cdot K_\text{eig})$近似获取前$K_\text{eig}$个特征对\cite{Shuman2013GSP}，是第~3~章算法~\ref{alg:ch3-lanczos}的预计算基础。图信号处理（GSP）框架\cite{Ortega2018GSP}揭示低频分量编码全局平滑信息、高频分量刻画局部差异的基本规律，为频率接口设计必要性提供了谱论背景；本文方法在频率维采用等权实例化以保持任务通用性，精细频率感知设计留待后续工作。

以Gilmer等人（2017）\cite{Gilmer2017MPNN}的消息传递框架（MPNN）来看，$L$层图神经网络的推理复杂度为$O(L\cdot m\cdot d)$；更重要的是，节点$i$的有效感受野在树状图上可达$O(\bar{d}^L)$量级，使在线单节点查询时延随层数快速放大。GCN\cite{Kipf2017GCN}、GAT\cite{Velickovic2018GAT}、GraphSAGE\cite{Hamilton2017GraphSAGE}、GIN\cite{Xu2019GIN}、SGC\cite{Wu2019SGC}等主流模型在不同程度上面临这一推理瓶颈，且层数增加还会引发过平滑（节点表示趋同）和过压缩（远程信息丢失）两类固有局限\cite{Li2018OverSmoothing,Alon2021OverSquashing}。前者通常与反复低通传播导致的表示收缩有关，后者则更多与图上的拓扑瓶颈和长程信息压缩有关；二者都可从图的谱性质与结构性质中获得分析线索，但具体机理并不相同。这说明单纯增加层数无法从根本上改善推理效率，稀疏化分解方向是有依据的技术选择。

% ---- 图~\ref{fig:ch2-bottleneck}~占位符 ----
\begin{figure}[htbp]
  \centering
    \includegraphics[width=1\linewidth]{2-2.pdf}
  % \setlength{\fboxsep}{1.2em}
  % \fbox{\parbox{0.88\textwidth}{\small\centering
  %   \textbf{【占位图~\ref{fig:ch2-bottleneck}~】}\\[0.6em]
  %   图题：消息传递GNN感受野扩张与在线推理瓶颈示意图。\\[0.4em]
  %   表达内容：图中左侧给出1层、2层、3层消息传递对应的感受野递增示意，右侧给出在线单节点查询时邻居收集与内存访问开销随层数增加而快速放大的趋势。横轴为层数$L$，纵轴为潜在访问节点规模；图中需对比“完整多层传播”和“稀疏近似推理”两类口径，突出相关工作中效率问题的来源。
  % }}
  \caption{消息传递GNN感受野扩张与在线推理瓶颈示意图}
  \label{fig:ch2-bottleneck}
\end{figure}

%%============================================================
\section{高效与稀疏图神经网络方法}
\label{sec:ch2-efficient}

针对感受野指数扩张瓶颈，学界发展出三条主要技术路线。基于采样的方法（GraphSAGE\cite{Hamilton2017GraphSAGE}、FastGCN\cite{Chen2018FastGCN}、GraphSAINT\cite{Zeng2020GraphSAINT}）通过控制每层邻居集合规模缓解训练阶段的邻域爆炸问题；在推理阶段，这类方法也可沿用采样近似控制开销，但若希望降低采样方差或保持预测稳定性，通常仍需访问较大邻域，因此对严格在线单节点查询时延的优化并不如稀疏分解路线彻底。传播与特征变换解耦的方法可进一步分为两类：一类以SGC\cite{Wu2019SGC}、SIGN\cite{Frasca2020SIGN}为代表，适合将多阶传播结果离线预计算后在推理期执行轻量变换；另一类以APPNP\cite{Gasteiger2019APPNP}为代表，通过个性化传播将预测与图扩散解耦，但其传播过程通常仍以迭代近似形式实现。总体而言，这类方法均显著缓解了图传播与在线推理开销的强耦合。

稀疏分解框架代表了更彻底的推理加速思路。Hu等提出的SDGNN\cite{Hu2024SDGNN}以稀疏矩阵分解逼近多层GNN输出，其推理阶段写为
\begin{equation}
  \hat{H} = \Theta^\top \tilde{\Phi},
  \label{eq:ch2-sdgnn-infer}
\end{equation}

其中$\tilde{\Phi}\in\mathbb{R}^{n\times d}$为教师GNN预计算的候选特征矩阵，稀疏权重矩阵$\Theta$（每列对应一个节点，最多$\bar{k}$个非零元）通过LARS算法离线求解以下Lasso问题：
\begin{equation}
  \min_{\theta_i} \|H_{i,:}^* - \tilde{\Phi}^\top\theta_i\|_2^2 + \lambda_{\mathrm{reg}}\|\theta_i\|_1, \quad i=1,\ldots,n.
  \label{eq:ch2-sdgnn-lasso}
\end{equation}

推理复杂度降至$O(n\bar{k}d)$，实现了与多层图传播的彻底解耦。然而，式~\eqref{eq:ch2-sdgnn-lasso}中的正则系数$\lambda_{\mathrm{reg}}$对全图所有节点统一施加相同的$\ell_1$惩罚，是SDGNN最突出的结构性约束：高度枢纽节点与边缘节点、同配区域与异配区域的节点在拓扑角色和频域参与方式上本应各异，全局均一系数无法自动适配这种差异。这一根本不足引出了第~3~章至第~5~章方法链中三维自适应稀疏化机制的核心研究动机。近期也出现了若干与图稀疏化、动态图注意或拉普拉斯稀疏化相关的工作。PPRGo（Bojchevski等，2020）\cite{Bojchevski2020PPRGo}基于个性化 PageRank 传播近似，以线性时间实现快速全图推理，在大规模图上推理效率较高，但其传播系数对全图共享，节点级自适应性不足。GLNN（Zhang等，2022）\cite{Zhang2022GLNN}通过知识蒸馏将 GNN 压缩为完全去除图依赖的 MLP，推理期无需访问图结构，达到极低延迟，但由此放弃了对图拓扑的利用。上述两类方法均未从节点谱域特征出发设计自适应稀疏化预算，与 SDGNN 稀疏分解路线和本文的关注点存在本质差异，本文在第~6~章实验中将其纳入效率对比参照，以说明不同技术路线的效率—精度权衡边界。就本文关心的问题而言，现有工作仍缺少一条从可观测图谱性质出发、再过渡到节点级稀疏预算设计准则的系统分析链。

%%============================================================
\section{复杂图结构与谱视角相关方法}
\label{sec:ch2-heterophily}

第~\ref{sec:ch2-efficient}~节揭示了SDGNN全局均一正则系数对节点异质性的结构性失配。这种异质性的深层来源，在于图上信号的频率分布差异。本节从复杂图结构下的频率响应出发，梳理现有相关方法，为频率接口设计必要性提供背景动机。

\subsection{图结构差异与频率响应}
\label{subsec:ch2-homophily}

已有研究普遍观察到，不同结构类型的图在有效判别频段上可能表现出明显差异：引文网络等较平滑图结构通常更依赖低频聚合，而网页/维基图等复杂结构数据常需要保留更多中高频成分\cite{Pei2020GeomGCN,Zhu2020H2GCN}。这类差异说明频率接口不宜简单等同于单一低通假设，但也不意味着必须引入标签统计量作为方法输入。本文据此仅保留“频率接口”这一设计位置，并在正文主线中采用等权实例化，使后续参数体系完全由纯结构统计量驱动。

\subsection{代表性复杂图结构处理方法}
\label{subsec:ch2-heteromethods}

从频率视角看，结构较平滑的图中判别信号更偏低频分量，标准低通聚合通常有效；而在结构更复杂的图中，判别信息往往更多分布于中高频，统一低通聚合会主动抑制有用信号。现有方法沿三条路线回应这一问题。H2GCN（Zhu等，NeurIPS 2020）\cite{Zhu2020H2GCN}分离ego与邻居表示并引入二阶邻域；GPR-GNN（Chien等，ICLR 2021）\cite{Chien2021GPRGNN}以可学习广义PageRank系数$\{\gamma_k\}$自适应调节传播权重，频率响应由数据驱动确定：
\begin{equation}
  H^{(K)}=\sum_{k=0}^K\gamma_k(I-\tilde{\mathcal{L}})^k H^{(0)};
  \label{eq:ch2-gprgnn}
\end{equation}

FAGCN（Bo等，AAAI 2021）\cite{Bo2021FAGCN}通过边级自门控系数$\alpha_{ij}^G\in[-1,1]$在消息传播中自适应整合低频与高频成分；LINKX\cite{Lim2021LINKX}和GloGNN\cite{Li2022GloGNN}则分别从结构—特征并行嵌入和全局节点关系聚合角度提供不同解决路径。

然而，上述方法的自适应粒度均不直接等同于“节点级稀疏预算”：GPR-GNN的$\{\gamma_k\}$对全图共享，FAGCN以边为粒度调节，二者均未形成“节点级局部频率环境$\rightarrow$节点级频率响应$\rightarrow$节点级稀疏预算”的完整映射；且这些方法推理期仍依赖完整图传播，与$O(n\bar{k}d)$稀疏分解接口不兼容。节点谱能量作为刻画节点在各频率分量上参与强度的关键指标，在上述方法中均未被显式引入。这三点共同界定了后续方法创新的边界，详见第~\ref{sec:ch2-summary}~节研究空白分析。上述频谱现象为频率接口设计必要性提供了背景动机；本文方法在频率维采用等权实例化以保持任务通用性，更精细的频率感知设计留待后续工作。

% ---- 图~\ref{fig:ch2-heterophily-spectrum}~占位符 ----
\begin{figure}[htbp]
  \centering
    \includegraphics[width=1\linewidth]{2-3.pdf}
  % \setlength{\fboxsep}{1.2em}
  % \fbox{\parbox{0.88\textwidth}{\small\centering
  %   \textbf{【占位图~\ref{fig:ch2-heterophily-spectrum}~】}\\[0.6em]
  %   图题：同配图与异配图的频率响应差异示意图。\\[0.4em]
  %   表达内容：图中左侧画出同配图标签信号偏低频的示意频谱，右侧画出异配图中判别信息向中高频扩展的示意频谱；中间用标准低通聚合、可学习多项式滤波和高低频自适应融合三类方法的响应曲线作对比，用以说明为何异配性处理方法会成为后续频率维设计的直接背景。
  % }}
  \caption{不同图结构的频率响应差异示意图}
  \label{fig:ch2-heterophily-spectrum}
\end{figure}

%%============================================================
\section{本章小结与研究空白}
\label{sec:ch2-summary}

前述三节系统梳理了与本文核心方法链直接相关的理论基础与代表性方法，由此可识别出现有工作在三个层面的不足，并与第~3~章至第~5~章方法设计逐一衔接。

就节点级自适应稀疏化的设计依据而言，SDGNN的全局统一$\lambda_{\mathrm{reg}}$（式~\eqref{eq:ch2-sdgnn-lasso}）与DropEdge\cite{Rong2020DropEdge}的随机删边均对全图所有节点施加相同稀疏惩罚，未能区分不同拓扑角色节点的信息整合需求；现有工作尚缺少一条从可观测谱性质出发推导节点级稀疏预算分配依据的分析路径，该不足由第~3~章通过问题形式化与四类图复杂度指标的建立完成可计算化铺垫，并在第~4~章进一步落实为频率维、节点维与跳距维的三维自适应参数设计。在节点级频率特性与稀疏化的联合设计层面，GPR-GNN的全局共享传播系数、FAGCN的边级门控、H2GCN的高阶邻域分离各自在一定程度上缓解了异配问题，却均未形成从节点局部频率特性到节点级稀疏预算的完整设计链；第~4~章在保持任务通用性的前提下以等权谱能量接口应对这一空白，将精细频率感知设计留作后续研究。此外，现有工作在节点级自适应稀疏化的参数解释、复杂度收益与误差关系上仍缺少统一讨论，本文在第~5~章对误差分析与复杂度分析进行系统整合，并与第~3~章、第~4~章的方法输出形成相互衔接的说明。

\begin{table}[htbp]
  \centering
  \small
  \setlength{\tabcolsep}{3pt}
  \begin{threeparttable}
    \caption{现有方法研究空白与本文对应章节方案汇总}
    \label{tab:ch2-research-gaps}
    \begin{tabular}{@{}>{\centering\arraybackslash}m{1.2cm}>{\centering\arraybackslash}m{3.2cm}>{\centering\arraybackslash}m{2.8cm}>{\centering\arraybackslash}m{2.4cm}>{\centering\arraybackslash}m{3.6cm}@{}}
      \toprule
      空白 & 核心局限描述 & 代表性相关方法 & 本文对应章节 & 核心技术手段 \\
      \midrule
      空白一 & 全局统一稀疏化，缺乏谱理论指导的节点级自适应框架
        & SDGNN（$\lambda_{\mathrm{reg}}$全局统一）、DropEdge & 第~3~章（第~4~章继续展开）
        & 问题形式化、指标提取与三维参数设计的前置接口\\[3pt]
      空白二 & 复杂结构图处理仍多采用全局频率权重，节点级频率特性与稀疏预算未建立联系
        & GPR-GNN、FAGCN、H2GCN & 第~4~章
        & 等权实例化频率接口，以节点谱能量等结构统计量驱动，保留频率感知扩展接口 \\[3pt]
      空白三 & 现有工作在节点级自适应稀疏化的参数解释、复杂度收益与误差关系上仍缺少统一讨论
        & 谱图理论、稀疏近似与复杂度分析（分散发展）
        & 第~5~章
        & 误差分析与适用条件说明、端到端复杂度分析\\
      \bottomrule
    \end{tabular}
    \begin{tablenotes}
      \small
      \item 注：本表中代表性方法的信息均引自本章正文各节；“核心技术手段”一列仅作方向性概述，不在本章提前给出后续章节的原创符号；相关正式记号统一在对应章节引言的符号约定表中定义。
    \end{tablenotes}
  \end{threeparttable}
\end{table}

\begin{figure}[htbp]
  \centering
    \includegraphics[width=1\linewidth]{2-4.pdf}
  % \setlength{\fboxsep}{1.2em}
  % \fbox{\parbox{0.88\textwidth}{\small\centering
  %   \textbf{【占位图~\ref{fig:ch2-research-gaps}~】}\\[0.6em]
  %   图题：三大研究空白与本文方法链的对应关系图。\\[0.4em]
  %   表达内容：图中左侧纵向排列三个矩形框，分别标注三大研究空白（“空白一：全局统一稀疏化/缺乏谱理论指导”、“空白二：节点级同配差异/联合设计缺失”、“空白三：谱稀疏化误差/缺乏统一分析框架”），右侧纵向排列三个模块框（“第~3~章：问题形式化与指标体系”、“第~4~章：三维自适应机制”、“第~5~章：算法整合与理论保证”），中间用带颜色的箭头标注对应关系：空白一$\rightarrow$第~3~章（标注“问题形式化与指标体系”），空白二$\rightarrow$第~4~章（标注“基于同配差异的频率维参数化”），空白三$\rightarrow$第~5~章（标注“谱保持误差、信息衰减界与端到端复杂度分析”），强调三大空白具有内在层级关联而非简单并列；图注说明各箭头标注仅为预告性说明，具体方法定义留待对应章节展开。
  % }}
  \caption{三大研究空白与本文方法链的对应关系图}
  \label{fig:ch2-research-gaps}
\end{figure}

以上三大空白由本文第~3~章至第~5~章的三段式方法链依次填补：第~3~章从问题形式化与指标体系出发，第~4~章实现三维自适应机制设计，第~5~章完成误差分析与复杂度分析说明。本章建立的符号体系将在后续章节中直接复用。具体地，$\tilde{\mathcal{L}}$、$\lambda_\text{gap}$、$m = \mathrm{nnz}(A) = 2|E|$ 及SDGNN的式~\eqref{eq:ch2-sdgnn-infer}--\eqref{eq:ch2-sdgnn-lasso}将作为第~3~章问题形式化与指标体系构建的直接起点（教师GNN目标输出 $H^*$ 在第~3~章第~3.1.2~节正式定义）；本章关于复杂图结构与频谱分析的讨论构成频率接口设计必要性的背景动机，第~4~章频率维采用等权实例化，以节点谱能量等纯结构统计量驱动，不引入额外的标签一致性统计量；谱稀疏化误差、信息衰减与复杂度分析之间的统一关系，则在第~5~章继续展开。
